%% kapa: exam=1 round=29 year=2018 period=1 form=A subject=civil_law subject_ko=민법 seq=1 number=1
%% source: Q-Net 감정평가사 기출문제 (cst00309) artlSeq=5202808 file=extracted/1008/5202808/2018년 제29회 감정평가사 1차시험 기출문제/2018년 제29회 감정평가사 1차시험 기출문제/1교시 A형.pdf
%% conversion: kordoc 4.24.51 -> tools/gen1.py (auto); review=proofread-sonnet; flags=none
\begin{question}{1}
부재와 실종에 관한 설명으로 옳지 않은 것은? (다툼이 있으면 판례에 따름)
\begin{choices}
  \item 부재자로부터 재산처분권을 위임받은 재산관리인은 그 재산을 처분함에 있어서 법원의 허가를 받을 필요가 없다.
  \item 제1순위 상속인이 있는 경우에 제2순위 상속인은 실종선고를 청구할 수 있는 이해관계인이 될 수 없다.
  \item 부재자 재산관리인의 재산처분행위를 허가하는 법원의 결정은 기왕의 처분행위를 추인하는 방법으로도 할 수 있다.
  \item 실종선고를 받은 자가 실종기간 동안 생존하였다는 사실이 밝혀진 경우, 실종선고의 취소 없이도 이미 개시된 상속을 부정할 수 있다.
  \item 피상속인의 사망 후 그 상속인에 대한 실종선고가 이루어졌으나 실종기간 만료 시점이 피상속인의 사망 이전인 경우, 실종선고된 자는 상속인이 될 수 없다.
\end{choices}
\end{question}
\answer{4}
